\subsection{\texorpdfstring{将 \(k[\mathbf{P}]\) 注入 \(\mathbf{S}(\mathfrak{h})^{*}\)}{将 k{[}\textbackslash mathbf\{P\}{]} 注入 \textbackslash mathbf\{S\}(\textbackslash mathfrak\{h\})\^{}\{*\}}}

映射 \(p \mapsto \exp p\) 从 P 到 \(\mathbf{S}(\mathfrak{h})^{*}\)，是从加法群 P 到赋予乘法结构的 \(\mathbf{S}(\mathfrak{h})^{*}\) 的同态（第 1 节）。因此，存在唯一的同态 \(\psi\)，从幺半群 P 的代数 k{[}P{]} 到代数 \(\mathbf{S}(\mathfrak{h})^{*}\)，使得

\[
\psi (e ^ {\lambda}) = \exp (\lambda) \quad (\lambda \in \mathrm{P})
\]

（采用 §7，第 7 节的记号）。根据第 1 节，\(\psi\) 是单射。通过结构迁移，有 \(\psi(w(e^{\lambda})) = w(\psi(e^{\lambda}))\)，对所有 \(\lambda \in \mathrm{P}\) 和所有 \(w \in \mathrm{W}\) 成立。因此，若 \(k[\mathrm{P}]^{\mathrm{W}}\)（分别为 \(\mathbf{S}(\mathfrak{h})^{*\mathrm{W}}\)）表示 \(k[\mathrm{P}]\)（分别为 \(\mathbf{S}(\mathfrak{h})^{*}\)）中在 \(\mathrm{W}\) 下不变的元素所成的集合，则 \(\psi(k[\mathrm{P}]^{\mathrm{W}}) \subset \mathbf{S}(\mathfrak{h})^{*\mathrm{W}}\)。

命题 1. 设 \(S^m(\mathfrak{h}^*)^{\mathrm{W}}\) 是 \(S^m(\mathfrak{h}^*)\) 中在 W 下不变的元素集合。则 \(\mathrm{pr}_m(\psi(k[\mathrm{P}]^{\mathrm{W}})) = \mathbf{S}^m(\mathfrak{h}^*)^{\mathrm{W}}\)。

由前述结果显然有 \(\mathrm{pr}_{m}(\psi(k[\mathrm{P}]^{\mathrm{W}})) \subset \mathbf{S}^{m}(\mathfrak{h}^{*})^{\mathrm{W}}\)。\(\mathbf{S}^{m}(\mathfrak{h}^{*})\) 的每个元素都是下列形式的元素的 k-线性组合：

\[
\mathrm{pr} _ {m} (\exp \lambda) = (\mathrm{pr} _ {m} \circ \psi) (e ^ {\lambda})
\]

其中 \(\lambda \in \mathrm{P}\)（引理 1）。因此，\(\mathbf{S}^m (\mathfrak{h}^*)^{\mathrm{W}}\) 的每个元素都是下列形式的元素的线性组合：

\[
\sum_ {w \in \mathrm{W}} w ((\mathrm{pr} _ {m} \circ \psi) (e ^ {\lambda})) = (\mathrm{pr} _ {m} \circ \psi) \left(\sum_ {w \in \mathrm{W}} w (e ^ {\lambda})\right),
\]

其中每一个都属于 \(\mathrm{pr}_m(\psi (k[\mathrm{P}]^{\mathrm{W}}))\)。

命题 2. 设 E 是有限维 g-模。设 U(h) = S(h) 是 h 的包络代数。若 \(u \in U(h)\)，则

\[
\operatorname{Tr} (u _ {\mathrm{E}}) = \langle \psi (\operatorname{ch} \mathrm{E}), u \rangle .
\]

只需处理 \(u = h_{1} \ldots h_{m}\) 且 \(h_{1}, \ldots, h_{m} \in h\) 的情形。对于所有 \(\lambda \in P\)，令 \(d_{\lambda} = \dim E^{\lambda}\)。则 \(\operatorname{ch} E = \sum_{\lambda} d_{\lambda} e^{\lambda}\)，所以 \(\psi(\operatorname{ch} E) = \sum_{\lambda} d_{\lambda} \exp(\lambda)\)，从而

\[
\begin{array}{r l} \langle \psi (\mathrm{chE}), u \rangle & = \sum_ {\lambda} d _ {\lambda} \langle \exp \lambda , h _ {1} \dots h _ {m} \rangle \\ & = \sum_ {\lambda} d _ {\lambda} \lambda (h _ {1}) \dots \lambda (h _ {m}) \quad (\text {no. 1}) \\ & = \operatorname{Tr} u _ {\mathrm{E}}. \end{array}
\]

推论 1. 设 \(\mathrm{U}(\mathfrak{g})\) 是 \(\mathfrak{g}\) 的包络代数。设同态 \(\zeta : \mathrm{U}(\mathfrak{g})^* \to \mathrm{U}(\mathfrak{h})^* = \mathbf{S}(\mathfrak{h})^*\) 是典则嵌入 \(\mathrm{U}(\mathfrak{h}) \to \mathrm{U}(\mathfrak{g})\) 的转置。下图交换

\[
\begin{array}{c c c} \mathcal {R} (\mathfrak {g}) & \stackrel {{\text {ch}}} {{\longrightarrow}} & \mathbf {Z} [ \mathrm{P} ] \\ \operatorname{Tr} \Big \downarrow & & \Big \downarrow_ {\psi} \\ \mathrm{U} (\mathfrak {g}) ^ {*} & \stackrel {{\zeta}} {{\longrightarrow}} & \mathbf {S} (\mathfrak {h}) ^ {*}. \end{array}
\]

这只是命题 2 的一种重新表述。

推论 2. 设 m 是满足 \(\geq 0\) 的整数。\(\mathbf{S}^{m}(\mathfrak{h}^{*})^{\mathrm{W}}\) 的每个元素都是 h 上形如 \(x \mapsto \operatorname{Tr}(\rho(x)^{m})\) 的多项式函数的线性组合，其中 \(\rho\) 是 g 的有限维线性表示。

根据命题 1，\(\mathbf{S}^m (\mathfrak{h}^*)^{\mathrm{W}} = (\mathrm{pr}_m\circ \psi)(k[\mathrm{P}]^{\mathrm{W}})\)。现在 \(\mathbf{Z}[\mathrm{P}]^{\mathrm{W}} = \operatorname {ch}\mathcal{R}(\mathfrak{g})\)（§7，第 7 节，定理 2 (ii)）。因此，根据第 VI 章，§3，第 4 节，引理 3，\(\psi (k[\mathrm{P}]^{\mathrm{W}})\) 是 \(k\)-向量子空间，即 \(\mathbf{S}(\mathfrak{h})^{*}\) 中由 \(\psi (\operatorname {ch}\mathcal{R}(\mathfrak{g})) = \zeta (\operatorname {Tr}\mathcal{R}(\mathfrak{g}))\) 生成的。所以，\(\mathbf{S}^m (\mathfrak{h}^*)^{\mathrm{W}}\) 是 \(\mathbf{S}^m (\mathfrak{h}^*)\) 中由 \((\mathrm{pr}_m\circ \zeta \circ \mathrm{Tr})(\mathcal{R}(\mathfrak{g}))\) 生成的向量子空间。但是，若 \(\rho\) 是 \(\mathfrak{g}\) 的有限维线性表示，

\[
((\operatorname{pr} _ {m} \circ \zeta \circ \operatorname{Tr}) (\rho)) (x) = \left\langle (\zeta \circ \operatorname{Tr}) (\rho), \frac {x ^ {m}}{m !} \right\rangle = \frac {1}{m !} \operatorname{Tr} (\rho (x) ^ {m})
\]

对所有 \(x \in h\) 都成立。

